\section{Transitional Effects: Multi-Point T-H Heater Integration (RQ1/RQ2)}
\label{sec:method_pinch}

Following supervisor feedback to model transitional thermophysical effects,
the two-state-point energy balance of the primary heater is replaced by a
nodal integration along the exchanger. The working fluid is discretised in
200 equal enthalpy steps between the heater inlet and the turbine inlet, and
the temperature and $c_p$ are evaluated from real-fluid properties at every
node, which resolves the pseudo-critical peak. The same procedure is applied
to the steam boiler (where CoolProp returns the saturation temperature
through the two-phase region) and to the sCO$_2$ heater, so that all three
cycles are integrated with the waste-gas stream by one method.

The waste gas is treated as an ideal gas with constant $c_p$ in counter-flow.
Its temperature therefore falls \emph{linearly with heat duty}, not with
position or with node index:
\begin{equation}
T_{gas}(q) = T_{gas,out} + \left(T_{src} - T_{gas,out}\right)\frac{q}{Q}
\label{eq:gas_profile}
\end{equation}
where $q$ is the cumulative heat transferred from the cold end and $Q$ the
total heater duty per kilogram of working fluid. The distinction matters: a
profile that is linear in node index implicitly assumes the gas gives up equal
heat per node, which is false wherever the working fluid's $c_p$ varies. The
Interim Report used an index-linear profile, which produced a spurious
constant 10\,K difference along the whole heater; the corrected profile is
used here.

Two constraints close the problem. The turbine inlet is set 10\,K below the
source temperature (hot-end approach), and the gas outlet temperature is
found by bisection so that the minimum local temperature difference equals
the 10\,K approach:
\begin{equation}
\mathrm{PPTD} = \min_{i}\!\left[T_{gas,i} - T_{wf,i}\right] = 10\,\text{K}
\label{eq:pptd}
\end{equation}
This is the maximum heat that the stream can give up to that cycle. Because
the heater inlet temperature of the working fluid is the floor for the gas
outlet, a cycle whose heater inlet is hot (a recuperated sCO$_2$ cycle, for
example) can recover less of the stream than one whose heater inlet is cold.
The conductance is then accumulated segment by segment,
\begin{equation}
UA = \sum_{i} \frac{\Delta Q_i}{\mathrm{LMTD}_i}, \qquad
\mathrm{LMTD}_i = \frac{\Delta T_{i} - \Delta T_{i+1}}{\ln(\Delta T_{i}/\Delta T_{i+1})}
\label{eq:seg_UA}
\end{equation}
which is exact for counter-flow with piecewise-linear profiles. The two-point
LMTD result, $UA_{2pt} = Q/\mathrm{LMTD}(\Delta T_{cold},\Delta T_{hot})$, is
computed alongside for comparison (Section~\ref{sec:rq2_hx_sizing}).

\section{Equipment Sizing and Capital Cost (RQ2)}
\label{sec:method_rq2}

\subsection{Site Design Points}
\label{sec:method_sitedesign}

Each Australian site (Table~\ref{tab:casestudies}) is a fixed waste-gas
stream with source temperature $T_{src}$ and mass flow $\dot{m}_{gas}$. Each
cycle is integrated with the stream through the heater model of
Section~\ref{sec:method_pinch}. For the CO$_2$ cycles the pressures follow the
RQ1 rules. For the steam cycle the boiler pressure is a free design variable:
the evaporation temperature fixes an internal pinch at the start of boiling,
so a higher pressure raises cycle efficiency but prevents the gas from being
cooled as far. The pressure is therefore chosen by a golden-section search to
maximise the net work per kilogram of waste gas,
\begin{equation}
w_{gas} = \eta_{th}\; c_{p,gas}\,\left(T_{src} - T_{gas,out}\right),
\qquad \dot{W}_{site} = \dot{m}_{gas}\, w_{gas}
\label{eq:wgas}
\end{equation}
with $c_{p,gas} = 1.10$\,kJ/(kg$\cdot$K). Equation~\eqref{eq:wgas} is the
quantity a plant owner actually receives from a given stream; it rewards
both conversion efficiency and depth of heat recovery, and it is used
throughout RQ2--RQ4. The recuperator of each CO$_2$ cycle is discretised in
the same way as the heater (200 enthalpy steps with real properties on both
sides) to obtain its $UA$ and to check that no internal temperature crossover
occurs. The heat-rejection exchanger (condenser or gas cooler) is assigned a
mean temperature difference of 10\,K to the cooling medium for all three
cycles.

\subsection{Turbomachinery Compactness}

Turbine size parameter and volumetric flow ratio are calculated from
Macchi and Astolfi~\cite{Macchi2017}:
\begin{equation}
\mathrm{SP} = \frac{\dot{V}_{out}^{0.5}}{\Delta h_{is}^{0.25}}, \qquad
\mathrm{VFR} = \frac{v_{out}}{v_{in}}
\label{eq:sp_vfr}
\end{equation}
where $\dot{V}_{out}$ is the volumetric flow rate at turbine exhaust (m$^3$/s
per kg/s of working fluid) and $\Delta h_{is}$ is the isentropic enthalpy
drop (kJ/kg). A VFR below approximately 5 is taken as the threshold for
single-stage radial turbine feasibility~\cite{Macchi2017}. Both are evaluated
at the site design points.

\subsection{Heat Exchanger Area}

For the illustrative comparison of sizing methods, the tCO$_2$ heater area is
reported as $A = UA/U$ with the overall coefficient computed from the PCHE
resistance network
\begin{equation}
\frac{1}{U} = \frac{1}{h_{CO_2}} + \frac{t_{wall}}{k_{wall}} + \frac{1}{h_{gas}}
\label{eq:U_overall}
\end{equation}
with $h_{CO_2} = 2000$\,W/m$^2$K (supercritical CO$_2$ side, conservative
PCHE estimate from Dostal~\cite{Dostal2004}), $h_{gas} = 500$\,W/m$^2$K
(exhaust gas forced convection), $t_{wall} = 2$\,mm, and $k_{wall} =
15$\,W/mK (SS316), giving $U = 380$\,W/m$^2$K. The comparison is made at
$\varepsilon = 0$ (pump exit to turbine inlet), which places the
pseudo-critical $c_p$ peak inside the heater and maximises the divergence
between the segmented and two-point results. Capital costs are computed from
$UA$ directly, so the value of $U$ does not affect the economics.

\subsection{Capital Cost Estimation}
\label{sec:method_cost}

Every cycle is costed with the same method and the same correlations, applied
component by component to the site design point scaled to the required net
output (Table~\ref{tab:cost_basis}). CO$_2$ turbines and all heat exchangers
use the NETL correlations of Weiland, Lance and Pidaparti~\cite{Weiland2019};
the steam turbine uses a power law fitted to the US~DOE small-turbine
benchmarks~\cite{DOE2012}; pumps and the sCO$_2$ compressor use Turton et
al.~\cite{Turton2012}. Turbine cost is evaluated on \emph{shaft} power
($\dot{W}_{net}/(1-\mathrm{BWR})$), which penalises the high back-work ratio
of the sCO$_2$ cycle correctly. The condenser or cooler, omitted at the
interim stage, is included for all three cycles because low-efficiency WHR
plants reject three to six units of heat per unit of electricity.

\begin{table}[H]
\centering
\caption{Capital cost basis. All correlations escalated to 2026 by CEPCI
(2001: 397; 2012: 584.6; 2017: 567.5; 2026: $\approx$830~\cite{CEPCI2026})
and converted at 1\,USD = 1.395\,AUD~\cite{RBA2026}; +35\,\% installation on
all equipment~\cite{Turton2012}.}
\label{tab:cost_basis}
\renewcommand{\arraystretch}{1.3}
\small
\begin{tabularx}{\textwidth}{>{\raggedright\arraybackslash}p{3.0cm}
  >{\raggedright\arraybackslash}p{5.2cm}
  >{\raggedright\arraybackslash}p{2.6cm}
  >{\raggedright\arraybackslash}X}
\toprule
\rowcolor{tableheader}
\textbf{Component} & \textbf{Correlation (cost year)} & \textbf{Scaling variable} & \textbf{Validity / note} \\
\midrule
\rowcolor{rowgray}
Primary heater, recuperator & $49.45\,UA^{0.7544}\,f_T$ (2017)~\cite{Weiland2019} & $UA$ (W/K) & $1.6\times10^5$--$2.1\times10^8$\,W/K; $f_T=1$ below 550\,\textdegree C \\
Condenser / cooler & $32.88\,UA^{0.75}$ (2017)~\cite{Weiland2019} & $UA$ (W/K) & 10\,K mean $\Delta T$ assumed \\
\rowcolor{rowgray}
CO$_2$ turbine & $406{,}200\,\dot{W}_{sh}^{0.8}$ (2017)~\cite{Weiland2019} & shaft MW & 8--35\,MW; extrapolated below \\
Steam turbine & $a\,\dot{W}_{sh}^{\,b}$, fitted to USD\,900/kW at 150\,kW and USD\,200/kW at 2\,MW (2012)~\cite{DOE2012} & shaft kW & back-pressure benchmark; condensing units cost more \\
\rowcolor{rowgray}
sCO$_2$ compressor & Turton centrifugal compressor (2001)~\cite{Turton2012} & kW & 450--3000\,kW \\
Feed / CO$_2$ pump & Turton centrifugal pump (2001)~\cite{Turton2012} & kW & 1--300\,kW; no pressure factor \\
\rowcolor{rowgray}
Generator & $108{,}900\,\dot{W}_e^{0.5463}$ (2017)~\cite{Weiland2019} & MWe & 4--750\,MWe; extrapolated below \\
\bottomrule
\end{tabularx}
\end{table}

The Interim Report used ad hoc power-law coefficients for the turbines and
approximate terminal temperature differences for the heaters; those were
replaced because (i)~the coefficients could not be traced to the cited
source and (ii)~the terminal-difference approximation produced unphysical
driving forces for the steam boiler and the sCO$_2$ heater. Specific capital
cost (AUD/kW) is reported at 100\,kW, 500\,kW, 1\,MW and 5\,MW net output at
each site's design conditions, and the cost breakdown at 1\,MWe.

\section{Techno-Economic Analysis (RQ3)}
\label{sec:method_rq3}

\subsection{Case Study Definitions}

Three Australian industrial case studies are drawn from ARENA~\cite{ARENA2023}
and A2EP~\cite{A2EP2022} data (Table~\ref{tab:casestudies}).

\begin{table}[H]
\centering
\caption{Australian industrial case study definitions.}
\label{tab:casestudies}
\renewcommand{\arraystretch}{1.3}
\begin{tabularx}{\textwidth}{>{\raggedright\arraybackslash}p{1.3cm}
  >{\raggedright\arraybackslash}p{3.9cm}
  >{\centering\arraybackslash}p{1.6cm}
  >{\centering\arraybackslash}p{1.5cm}
  >{\centering\arraybackslash}p{1.9cm}
  >{\raggedright\arraybackslash}X}
\toprule
\rowcolor{tableheader}
\textbf{Case} & \textbf{Industry / Site} & \textbf{$T_{src}$ (\textdegree C)} &
\textbf{$\dot{m}_{gas}$ (kg/s)} & \textbf{Price (AUD/MWh)} & \textbf{Rationale} \\
\midrule
\rowcolor{rowgray}
CS2 & Gas processing intercooler (Pilbara, WA) & 200 & 30 & 130 &
  Representative LNG train; isolated-grid price premium~\cite{AEMO2024} \\
CS1 & Cement preheater exhaust (Gladstone, QLD) & 350 & 50 & 90 &
  Dominant WHR opportunity in Australian cement \\
\rowcolor{rowgray}
CS3 & Electric arc furnace flue gas (Sydney, NSW) & 500 & 20 & 100 &
  High-temperature WHR in NEM-connected steelmaking \\
\bottomrule
\end{tabularx}
\end{table}

\subsection{Fixed-Source Basis}

The comparison is made on a \emph{fixed-source} basis: each cycle is sized to
the net output it can recover from the site's stream
(Eq.~\ref{eq:wgas}), and its capital cost follows from the cost model at that
size. This is the decision an operator faces --- the stream is given, the
plant is not. A fixed-output comparison (all cycles at 1\,MWe) is reported in
RQ2 for the capital cost; on that basis, because waste heat is free and the
plant life and capacity factor are common, LCOE is simply proportional to the
specific capital cost, $\mathrm{LCOE} = (\mathrm{CRF} + 0.025)\,
\mathrm{TEC}/(8760\,\mathrm{CF})$, and adds no information beyond RQ2.

\subsection{LCOE, NPV, Payback and IRR}

Levelised Cost of Electricity is calculated following Short et
al.~\cite{Short1995}:
\begin{equation}
\mathrm{LCOE} = \frac{C_{cap} \cdot \mathrm{CRF} + C_{OM}}{E_{gen}}
\label{eq:lcoe}
\end{equation}
where $C_{cap}$ is the total erected cost (AUD), $C_{OM}$ is annual
operations and maintenance (2.5\,\% of capital per year, consistent with WHR
industry practice~\cite{Baglietto2023}), $E_{gen}$ is annual electricity
generation (MWh/yr at 85\,\% capacity factor), and CRF is the capital
recovery factor at a 7\,\% real discount rate over a 20-year project life:
\begin{equation}
\mathrm{CRF} = \frac{r(1+r)^n}{(1+r)^n - 1} = 0.0944
\label{eq:crf}
\end{equation}
Net present value is the sum of discounted annual net revenue (electricity
revenue minus O\&M) over the project life less the initial capital. Payback
is reported as \emph{discounted} payback --- the year in which cumulative
discounted net cash flow first exceeds the capital, interpolated within the
year --- so that it is consistent with NPV and LCOE; the simple payback is
also recorded. Internal rate of return is the discount rate at which NPV is
zero. Electricity prices are drawn from AEMO spot market data and WA
wholesale tariffs~\cite{AEMO2024} (Table~\ref{tab:casestudies}). Waste heat
is free, no carbon price is applied, and the cost of mechanical chilling for
the tCO$_2$ condenser is not included (Section~\ref{sec:plow_sensitivity}
quantifies it separately).

\subsection{Monte Carlo Uncertainty Analysis}

Uncertainty is quantified by 10,000-sample Monte Carlo, varying three
parameters simultaneously with uniform distributions: capital cost
($\pm 30$\,\%, matching the stated uncertainty of the cost correlations),
electricity price ($\pm 40$\,\%) and capacity factor (60--95\,\%). The
primary output is $P(\mathrm{NPV} > 0)$, with P10/P50/P90 values of NPV and
LCOE. A breakeven analysis also reports the reduction in heat exchanger cost
(primary heater plus recuperator) at which each CO$_2$ cycle would match the
best cycle's LCOE at the same site.

\section{Off-Design Analysis (RQ4)}
\label{sec:method_rq4}

Part-load performance is evaluated at five load levels --- 100, 80, 60, 40,
and 20\,\% of design mass flow --- using a simplified Stodola efficiency
correction:
\begin{equation}
\eta_{od} = \eta_0 \cdot [1 - 0.3(1-f)^2]
\label{eq:stodola}
\end{equation}
where $f$ is the part-load mass-flow fraction and the parabolic coefficient
$0.3$ is calibrated to Yang et al.'s~\cite{Yang2024} tCO$_2$ part-load data.
The correction is applied to the turbine of every cycle and additionally to
the sCO$_2$ compressor. Sliding-pressure control is applied: $P_{high} =
P_{design}\,f$ for the SRC, $P_{high} = 20\,(0.5 + 0.5f)$\,MPa for sCO$_2$
and $P_{high} = P_{design}\,(0.6 + 0.4f)$ for tCO$_2$, with low-side
pressures and turbine inlet temperature held at their site design values.
The heater integration is not re-solved at part load.

Operational flexibility is scored using the Marchionni et
al.~\cite{Marchionni2020} framework with four criteria: turndown ratio (weight
0.30), efficiency retention at 60\,\% load (0.30), efficiency retention at
40\,\% load (0.25), and startup complexity (0.15). Retention values are
averaged over the three sites. Startup complexity scores are assigned
qualitatively based on the literature: SRC 0.55 (steam drum pressurisation,
blow-down, long warm-up required); sCO$_2$ 0.75 (no phase change, faster
startup, critical control near compressor operating point); tCO$_2$ 0.80
(simplest startup, liquid pump, no near-critical compressor
management~\cite{Marchionni2020}).

\section{Implementation and Reproducibility}
\label{sec:method_code}

All models are implemented in a single Python package (\texttt{co2whr}) with
one script per research question and a driver (\texttt{run\_all.py}) that
regenerates every figure and every number in this thesis in under two minutes
(Appendix~\ref{app:code}). Results are written to a machine-readable file
(\texttt{results.json}) from which the tables in Chapters~\ref{ch:results_rq1}
to~\ref{ch:results_rq4} were transcribed.
